\section{软商品与经济作物：白糖、棉花、苹果、红枣、花生、天然橡胶}
\label{sec:soft}

本章的七个品种（白糖、棉花、苹果、红枣、花生、天然橡胶、纸浆）与第 \ref{sec:grain} 章的
粮食油料有本质区别：\hl{它们的周期几乎不由中国的需求决定，而由产地气候、全球供需与
政策制度决定}。这决定了它们的周期性形态——不是猪周期那样的“产能$\to$价格”闭环，
而是\emph{外生冲击主导的脉冲式周期}。

\subsection{三类周期来源}
\label{sec:soft-drivers}

\begin{center}
\begin{tikzpicture}[font=\footnotesize,
  b/.style={rectangle, rounded corners=3pt, draw=AgriGreen, line width=0.9pt,
            fill=AgriPale, align=left, inner sep=5pt, text width=4.25cm,
            minimum height=2.6cm},
  h/.style={font=\small\heiti, AgriGreen, anchor=west}]
  \node[b] (t1) at (0,0)
    {\textbf{① 政策定价型}\\[2pt] 白糖、棉花\\[2pt]
     \scriptsize 目标价格补贴、进口配额/滑准税、国储与地方储备共同构成价格区间；
     周期表现为“政策调整$\to$种植面积变化$\to$产量变化”的 3$\sim$4 年节奏，
     价格振幅被制度性压缩。};
  \node[b, right=0.5cm of t1] (t2)
    {\textbf{② 产地气候型}\\[2pt] 苹果、红枣、花生\\[2pt]
     \scriptsize 供给高度集中于单一产区（新疆、陕西、山东），
     花期霜冻、成熟期连阴雨与“大小年”直接决定产量；
     周期表现为“减产$\to$价格暴涨$\to$补种/大年$\to$价格回落”的脉冲。};
  \node[b, right=0.5cm of t2] (t3)
    {\textbf{③ 全球供需型}\\[2pt] 天然橡胶、纸浆\\[2pt]
     \scriptsize 供给在东南亚/全球浆厂、需求在轮胎与造纸；
     中国的角色是消费方，价格由全球产能周期（胶树 7$\sim$10 年生长周期、浆厂投产节奏）
     与国际报价决定。};
  \node[h] at (-1.9,1.7) {周期来源三分法};
\end{tikzpicture}
\end{center}

\subsection{白糖：全球过剩与短缺的切换}
\label{sec:soft-sugar}

白糖的国内供给可以精确到榨季：2023/24 榨季产糖 996 万吨，
2024/25 榨季 \hl{1116.21 万吨（同比 +12.03\%）}，实现连续三个榨季的恢复性增长\cite{sugarassoc2025}；
2025/26 榨季最终统计产量 \hl{1327.62 万吨（同比 +16.68\%）}，
为历史第三高位\cite{yuntang2026sugar}。
需要说明的是，\hl{该榨季的产量口径在机构间存在显著分歧}：
中国糖业协会在 2025 年 11 月的预测为 \num{1170} 万吨，
而云糖网 2026 年 9 月的最终统计为 \num{1327.62} 万吨，两者相差 \num{157.6} 万吨
——\emph{榨季产量数据在榨季结束前具有很大的修正空间}，
这也是白糖价格在榨季中期出现预期修正式波动的原因之一。产量的连续回升来自糖料种植面积的恢复，
而种植面积的恢复又直接源于国内糖价在此前的上行（蔗农收益改善）——
这是政策定价型品种的典型传导：\emph{价格信号通过补贴与收购价传导到面积，再经过一个榨季传导到产量}。

需求侧的边际变化在国际市场：2025 年 1---10 月累计进口食糖 390.16 万吨（同比 +13.69\%），
糖浆与白砂糖预混粉进口约 100 万吨\cite{yuntang2026sugar}。
2025/26 年度全球食糖过剩量被国际糖业组织（ISO）从 220 万吨下调至 110 万吨，
并预计 2026/27 年度转为 \hl{20 万吨供应缺口}\cite{moa2026sugar}；
但机构间分歧极大：仅 2025/26 年度的过剩量估计，
就从 ISO 的 \num{110} 万吨（下调后）到 USDA 的 \num{1140} 万吨不等。
2026 年 8 月纽约原糖累计上涨超过 20\%、创 2010 年以来最大月度涨幅，
9 月 10 日 ICE 11 号原糖 10 月合约收于 18.76 美分/磅\cite{yuntang2026sugar}。

价格数据对这种“外强内弱”的结构给出了直接印证：2021 年 2 月至 2026 年 9 月，
郑商所白糖期货仅上涨 \hl{0.8\%}（图 \ref{fig:fut-norm-soft}），
而同期国际原糖经历了完整的“过剩$\to$短缺”切换。
国内价格被进口配额、国储与加工糖成本三重约束在一个窄区间内，
这正是白糖主窗口年化波动率仅 8.7\%、为全部农业期货中第二低的原因（图 \ref{fig:vol-bar}）。

\subsection{棉花：目标价格补贴与全球纺织需求}
\label{sec:soft-cotton}

棉花是中国农业中\hl{政策介入最深}的品种。2025 年全国棉花产量 664.1 万吨
（同比 +7.7\%），播种面积 2979.2 千公顷（+5.0\%），其中新疆产量 616.5 万吨
（占全国的 92.8\%）\cite{stats2025cotton}。供给的绝对集中意味着
\emph{新疆一地的天气就是全国供给的天气}。

政策方面，2026---2028 年新一轮新疆棉花目标价格政策延续每吨 18600 元的价格水平、
固定产量 510 万吨补贴，但有两处重要调整：\hl{质量补贴资金分配比例由 5\% 左右提高到 10\% 左右}，
\hl{长绒棉交售量补贴标准由细绒棉的 1.5 倍提高到 2 倍}\cite{xj2026cottonpolicy}。
政策取向从“保产量”转向“提质”——这与中央一号文件的种业振兴主线一致\cite{cccp2025doc1}。
进口方面，2025/26 年度累计进口棉花 158.5 万吨（同比 +50.7\%），
有关部门在 2025 年 8 月与 2026 年 3 月两次增发滑准税进口配额合计 50 万吨\cite{jinqiao2026cotton}。
\emph{配额的发放节奏本身构成一个政策周期}：配额宽松$\to$进口增加$\to$国内价格承压$\to$下一轮配额收紧。

从价格看，棉花期货在 2021---2026 年间上涨 3.5\%，但路径剧烈——先因 2021---2022 年
全球补库与新疆减产大涨至 21507 元/吨（主窗口最高），随后在纺织需求走弱中回落。
第 \ref{sec:corr-cycle} 节已经指出，棉花的周期分量与生猪的周期分量呈\hl{显著负相关（$-0.59$）}，
其根源正是两者驱动因素的完全错位：棉花由上一年度的全球纺织需求与新疆天气决定，
生猪由国内产能决定。

\subsection{苹果与红枣：真正的“产量周期”}
\label{sec:soft-fruit}

如果要在农业中寻找教科书式的“蛛网周期”，答案不是生猪，而是\hl{苹果}。

苹果的产量存在明确的\hl{“大小年”（隔年结果）}规律：连续两年产量相差 20\% 以上即认为出现
大小年现象，晚熟品种重于早中熟品种、老年树重于幼年树\cite{cma2015fruit}。
其机理是果树体内源激素与营养分配的自我调节——大年结果过多会消耗花芽分化所需的养分，
使次年成为小年。叠加花期霜冻与成熟期降水，形成“产量$\to$价格$\to$管理投入$\to$产量”的
短周期闭环。

数据清晰地展示了这一脉冲：全国苹果产量 2023/24 产季 3554.11 万吨\cite{fzzq2026apple}、
2024/25 产季升至 3734.97 万吨（+5.09\%）、
2025/26 产季回落至 \hl{3432.44 万吨（同比 -312 万吨，近七年最低）}\cite{fzzq2026apple}；
库存同步收缩，2025/26 产季全国入库峰值 735.77 万吨，为近五年最低、同比下降 104.57 万吨
\cite{fzzq2026apple}。产量与库存的双降推动苹果期货在 2026 年从低位反弹，
使其成为主窗口内唯一“逆势上涨”的种植类品种（+19.8\%，图 \ref{fig:fut-norm-soft}）。

红枣与苹果同属“产地气候 + 库存”驱动，但状态相反：
2024 产季新疆灰枣产量 75 万吨（较 2023 产季大增 67\%）\cite{mysteel2025jujube,qhrb2026jujube}，
2025 产季回落至 57.38 万吨（-23.49\%），但\emph{减产幅度有限且旧季库存高企}，
测算 2025/2026 产季结转库存将达 30 万$\sim$35 万吨、创近五年新高\cite{qhrb2026jujube}。
因此红枣呈现“减产不涨价”的形态：期货价格从主窗口初期的 15799 元/吨跌至 7579.7 元/吨
（\hl{-25.9\%}），主窗口年化波动率 25.9\%、仅次于生猪。
\hl{红枣是“库存周期压倒产量周期”的典型案例}——这一区分对第二部分的加工类企业尤其重要。

\subsection{花生：四年一轮回}
\label{sec:soft-peanut}

花生提供了一个“小品种、大波动”的完整样本。其需求结构为压榨 55\%$\sim$60\%、
食用 30\%$\sim$35\%、饲用 120 万$\sim$150 万吨，呈现“压榨定大盘、食用定品质”的双轨特征
\cite{nqw2026peanut}。由于压榨需求与油脂价格挂钩，花生在定价逻辑上
\hl{与豆油、棕榈油等油脂品种的联系强于与粮食的联系}。

供给端经历了一轮完整的四年周期：2022 年的短缺（价格高点）刺激种植面积扩张，
2024 年全国花生种植面积约 7256 万亩、产量 1898.3 万吨（连续第二年增产），
2025 年面积再增 4.01\% 至 457.49 万公顷（约 6862 万亩），为连增第三年\cite{nqw2026peanut}。
价格随之走弱，2025 年现货运行区间约 3.8---4.5 元/斤\cite{nqw2026peanut}。
期货端，花生主力合约从 2021 年的 10882.7 元/吨高位回落 24.2\%，
但主窗口年化波动率仅 12.1\%，是农产品期货中较低者——\emph{因为产量分散、
无单一产区天气风险，且价格由压榨利润锚定}。

\subsection{天然橡胶与纸浆：需求在外、供给在外}
\label{sec:soft-rubber}

这两个品种严格说来不属于“中国农业的周期”，而是\hl{中国农业资本参与全球资源品周期}的样本，
但它们在申万分类中对应种植业（天然橡胶）与轻工制造上游，且是农业上市公司
（如海南橡胶、部分林浆纸企业）的核心价格变量，因此纳入本章。

天然橡胶的全球供给已经连续多年偏紧：国际橡胶生产国协会（ANRPC）预计 2025 年全球天然橡胶
供应缺口将\hl{连续第五年存在}，全球产量仅增 0.3\% 至 1490 万吨、消费量增 1.8\% 至 1560 万吨；
印尼因农民转向油棕，2025 年产量预计同比下降 9.8\% 至 204 万吨，
越南减少 1.3\% 至 128 万吨\cite{anrpc2025rubber}。需求端，中国约 70\% 的天然橡胶用于轮胎生产
\cite{anrpc2025rubber}，因此橡胶周期本质上是\emph{汽车与物流周期}。
2021 年 2 月至 2026 年 9 月，天然橡胶上涨 \hl{24.8\%}、20 号胶上涨 43.3\%，
是主窗口内表现最好的农业相关品种之一。

纸浆的路径相反：2024 年主力合约波动区间 5500---6500 元/吨，2025 年四季度在进口浆报价
推动下反弹（12 月 31 日进口阔叶浆均价 4627.63 元/吨、较 9 月末 +9.08\%），
但 2026 年上半年转为偏弱下行，6 月末山东银星针叶浆现货 4720 元/吨（较年初 -880 元/吨），
港口库存维持 232 万吨以上的高位\cite{guoyuan2026pulp}。
结果是从 2021 年 2 月的高点回落 \hl{-29.8\%}，为本章跌幅最大的品种。
\hl{纸浆的周期是“浆厂投产 + 港口库存”的库存周期，与中国农产品供需无关}——
它在相关矩阵中相对孤立（图 \ref{fig:corr-heat}），正是这一独立性的统计表现。

\subsection{软商品的整体图景}
\label{sec:soft-summary}

\begin{figure}[htbp]
\centering
\begin{tikzpicture}
\begin{axis}[agri axis, yearticks,
  width=14.2cm, height=6.8cm,
  xlabel={年份（2021 年 2 月 $=100$）}, ylabel={归一化价格指数},
  xmin=2021.0, xmax=2026.9, ymin=55, ymax=165,
  xtick={2021,2022,2023,2024,2025,2026},
  legend style={at={(0.99,0.03)}, anchor=south east, legend columns=3,
                cells={anchor=west}, font=\tiny},
  legend entries={白糖, 棉花, 棉纱, 天然橡胶, 20 号胶, 苹果, 红枣, 纸浆, 花生},
]
\addplot[SoftRed, thick] table[x=x, y=SR]{data/clean/fut_norm_soft.dat};
\addplot[AquaTeal, thick] table[x=x, y=CF]{data/clean/fut_norm_soft.dat};
\addplot[AgriLight, thick] table[x=x, y=CY]{data/clean/fut_norm_soft.dat};
\addplot[InkGray, thick] table[x=x, y=RU]{data/clean/fut_norm_soft.dat};
\addplot[OilBlue, thick] table[x=x, y=NR]{data/clean/fut_norm_soft.dat};
\addplot[AgriGreen, thick] table[x=x, y=AP]{data/clean/fut_norm_soft.dat};
\addplot[AgriGold, thick] table[x=x, y=CJ]{data/clean/fut_norm_soft.dat};
\addplot[violet, thick] table[x=x, y=SP]{data/clean/fut_norm_soft.dat};
\addplot[Grain, thick] table[x=x, y=PK]{data/clean/fut_norm_soft.dat};
\node[font=\scriptsize, OilBlue, anchor=west] at (axis cs:2026.02,146) {20 号胶 143.3};
\node[font=\scriptsize, AgriGold, anchor=west] at (axis cs:2026.02,71) {红枣 74.1};
\node[font=\scriptsize, violet, anchor=west] at (axis cs:2026.02,66) {纸浆 70.2};
\end{axis}
\end{tikzpicture}
\caption{软商品与经济作物期货归一化走势（2021 年 2 月 $=100$）}
\label{fig:fut-norm-soft}
\srcfile{新浪财经期货主力连续合约；作者计算。}
\end{figure}

把本章与第 \ref{sec:grain} 章放在一起，可以得到软商品的三个总结论：

\begin{keybox}[软商品周期的三条规律]
\normalsize
\textbf{① 政策越深，波动越低。}白糖（8.7\%）、棉花（17.0\%）在配额、补贴与储备制度的
约束下波动受限；而苹果（20.4\%）、红枣（25.9\%）几乎无政策干预，波动完全由产量与库存释放。
\par\vspace{3pt}
\textbf{② 单一产区=单一风险。}新疆棉花（占全国 92.8\%）、新疆灰枣、
陕西与山东苹果（合计占全国一半以上），使这些品种的周期性表现为
“一地天气决定全国价格”的脉冲式波动，而非平滑周期。
\par\vspace{3pt}
\textbf{③ 库存决定“减产是否涨价”。}红枣在 2025 产季减产 23.5\% 的情况下价格仍下跌，
原因是结转库存创五年新高；苹果在减产 8.4\% 的同时库存降至五年最低，价格因此反弹。
\hl{产量周期必须与库存周期同时判断，这是种植类品种与养殖类品种最大的方法论差异}——
生猪无法库存，而果实可以。
\end{keybox}
